
\documentclass{article}
\usepackage{colortbl}
\usepackage{makecell}
\usepackage{multirow}
\usepackage{supertabular}

\begin{document}

\newcounter{utterance}

\centering \large Interaction Transcript for game `cladder', experiment `full\_v1.5\_default', episode 9908 with score{-}Qwen3.5{-}9B{-}sp.
\vspace{24pt}

{ \footnotesize  \setcounter{utterance}{1}
\setlength{\tabcolsep}{0pt}
\begin{supertabular}{c@{$\;$}|p{.15\linewidth}@{}p{.15\linewidth}p{.15\linewidth}p{.15\linewidth}p{.15\linewidth}p{.15\linewidth}}
    \# & \multicolumn{2}{c}{Player} && \multicolumn{2}{c}{Game Master} \\
    \hline

    \theutterance \stepcounter{utterance}  
    & & & \multicolumn{4}{p{0.6\linewidth}}{
        \cellcolor[rgb]{0.9,0.9,0.9}{
            \makecell[{{p{\linewidth}}}]{
                \texttt{\tiny{[P1$\langle$GM]}}
                \texttt{You are an expert in causal inference. The following question is not a typical commonsense query, but rather a meticulously designed question created by a professor specializing in causal inference, intended to assess the students' mastery of the course content.} \\
\texttt{Imagine a self{-}contained, hypothetical world with only the following conditions, and without any unmentioned factors or causal relationships: Zuph has a direct effect on wibl. Wibl has a direct effect on uvzi. We know that zuph causes not wibl, and we know that wibl causes not uvzi. Would an individual is not uvzi if zuph instead of not zuph?} \\
\texttt{Start your answer with 'Yes' or 'No', followed by additional reasoning or evidence to support your explanation.} \\
            }
        }
    }
     \\ \\

    \theutterance \stepcounter{utterance}  
    & \multicolumn{4}{p{0.6\linewidth}}{
        \cellcolor[rgb]{0.561,0.737,0.561}{
            \makecell[{{p{\linewidth}}}]{
                \texttt{\tiny{[P1$\rangle$GM]}}
                \texttt{No} \\
\\ 
\texttt{**Reasoning:**} \\
\\ 
\texttt{In this hypothetical world, we must strictly adhere to the provided causal structure and the specific truth values of the relationships.} \\
\\ 
\texttt{1.  **Causal Structure**: The problem states that \$Zuph \textbackslash{}rightarrow Wibl\$ and \$Wibl \textbackslash{}rightarrow Uvzi\$. This forms a linear chain: \$Zuph \textbackslash{}rightarrow Wibl \textbackslash{}rightarrow Uvzi\$.} \\
\texttt{2.  **Causal Directionality**: The problem explicitly states that "\$Zuph\$ causes **not** \$Wibl\$." In causal inference terms, this means the intervention \$do(Zuph = \textbackslash{}text\{true\})\$ results in the state \$Wibl = \textbackslash{}text\{false\}\$ (or "not wibl"). Similarly, "\$Wibl\$ causes **not** \$Uvzi\$" means that if \$Wibl\$ is true, \$Uvzi\$ becomes false. Conversely, if \$Wibl\$ is false, the direct causal push toward "not uvzi" from \$Wibl\$ is absent, but we must look at the baseline or the specific logic of the "not" operator.} \\
\texttt{    *   Let's analyze the logic of the "not" statements carefully.} \\
\texttt{    *   Statement A: \$Zuph \textbackslash{}implies \textbackslash{}neg Wibl\$.} \\
\texttt{    *   Statement B: \$Wibl \textbackslash{}implies \textbackslash{}neg Uvzi\$.} \\
\texttt{    *   Implicitly, in a binary world (True/} \\
            }
        }
    }
    & & \\ \\

    \theutterance \stepcounter{utterance}  
    & & & \multicolumn{2}{p{0.3\linewidth}}{
        \cellcolor[rgb]{0.9,0.9,0.9}{
            \makecell[{{p{\linewidth}}}]{
                \texttt{\tiny{[GM$|$GM]}}
                \texttt{No} \\
            }
        }
    }
    & & \\ \\

    \theutterance \stepcounter{utterance}  
    & & & \multicolumn{2}{p{0.3\linewidth}}{
        \cellcolor[rgb]{0.9,0.9,0.9}{
            \makecell[{{p{\linewidth}}}]{
                \texttt{\tiny{[GM$|$GM]}}
                \texttt{no} \\
            }
        }
    }
    & & \\ \\

    \theutterance \stepcounter{utterance}  
    & & & \multicolumn{2}{p{0.3\linewidth}}{
        \cellcolor[rgb]{0.9,0.9,0.9}{
            \makecell[{{p{\linewidth}}}]{
                \texttt{\tiny{[GM$|$GM]}}
                \texttt{game\_result = WIN} \\
            }
        }
    }
    & & \\ \\

\end{supertabular}
}

\end{document}
